\chapter{Точная LCF-проверка полярной оси межсекторной ковариации}
\label{ch:v8-cross-arrow-covariance-lcf-migration-gate}

\section{Точный алгебраический полярный подъём}

Численная полярная коизометрия исходного гейта восстановлена в поле
\begin{equation}
 \mathbb Q\!\left(\sqrt2,2\cos\frac\pi7\right).
\end{equation}
На физической матрице $A_0:\mathbb C^{11}\to\mathbb C^{10}$ построено
$U=(A_0A_0^*)^{-1/2}A_0$, причём равенство $UU^*=I_{10}$ проверено внутри
этого поля с нулевым остатком.

Полный линейный гессиан относительной кривизны вычислен без чисел с
плавающей точкой. После физической перестановки двенадцатимерный межсекторный блок
равен
\begin{equation}
 H_{\rm link}^{qX}=I_6\otimes B,
 \label{eq:v8-lcf-cross-six-pairs}
\end{equation}
а его угол с оставшимися пятнадцатью направлениями тождественно нулевой.

\section{Точная матрица пары}

Положим $c=\cos(\pi/7)$. Тогда
\begin{equation}
 B=\begin{pmatrix}
 \dfrac{124}{7}-4\sqrt2-\dfrac{96}{7}c^2&
 -\dfrac67+\dfrac{32}{7}c-\dfrac{32}{7}c^2\\[2mm]
 -\dfrac67+\dfrac{32}{7}c-\dfrac{32}{7}c^2&
 \dfrac{82}{7}-\dfrac{96}{7}c^2
 \end{pmatrix}.
 \label{eq:v8-lcf-exact-cross-pair}
\end{equation}
Точный критерий Сильвестра даёт положительную определённость, а положительный
дискриминант --- два различных собственных значения. Их десятичные образы
равны $0.2724480819\ldots$ и $1.2342661003\ldots$.

\section{Что фиксируется и что остаётся свободным}

Вакуумный парный гессиан имеет вид
\begin{equation}
 H_\eta^{qX}=\frac{32}{5}I_2+2\eta B.
\end{equation}
Поэтому при каждом $\eta>0$ он коммутирует с $B$ и имеет те же собственные
оси. Мягкая ковариационная ось равна
$55.4509155208\ldots^\circ$ и не зависит от $\eta$.

Однако отношение собственных дисперсий является нетривиальной функцией
$\eta$. Точно ненулевые производные также сохраняются у классической
скорости по коэффициенту действия, у гармонической скорости по
кинетическому множителю и у тепловой скорости по времени корреляции.
Следовательно, направление выведено, но форма эллипса и его масштаб нет.

\begin{theorem}[Точная полярная межсекторная ось]
Относительная полярная кривизна организует двенадцать межсекторных
направлений в шесть тождественных положительных пар, точно отделённых от
остальных пятнадцати мод. Все положительные $H_\eta$ выбирают одну общую
ось `QLYR--XLdR`.
\end{theorem}

\begin{nogocriterion}[Ось не является полной ковариацией]
Нельзя объявлять угол межсекторных типов углом CKM/PMNS или превращать его в
единственную скорость Крауса. Анизотропия зависит от свободного $\eta$, а
масштаб --- от коэффициента действия, кинетической нормы либо времени
корреляции.
\end{nogocriterion}

\begin{protocol}
\begin{enumerate}
 \item точное поле коэффициентов: \textbf{степень $6$};
 \item полярная коизометрия: \textbf{точно};
 \item одинаковых межсекторных пар: \textbf{$6$};
 \item связь с другими 15 модами: \textbf{ноль};
 \item парная матрица: \textbf{положительно определена};
 \item собственные значения: \textbf{различны};
 \item общая ось при $\eta>0$: \textbf{выведена};
 \item угол оси: \textbf{$55.4509155208\ldots^\circ$};
 \item анизотропия: \textbf{зависит от $\eta$};
 \item единственный масштаб и скорость: \textbf{не выведены}.
\end{enumerate}
\end{protocol}

Машинный сертификат сохранён в
\path{s2t_v8_cross_arrow_covariance_lcf_migration_gate_results.json}.